\pdfoutput=1
% !TEX program = xelatex
\documentclass[12pt,a4paper]{article}
% ================================================================
% Packages
% ================================================================
\usepackage{amsmath,amssymb,amsthm}
\usepackage{mathtools}
\usepackage{bm}
\usepackage{graphicx}
\usepackage{hyperref}
\usepackage{geometry}
\usepackage{enumitem}
\usepackage{booktabs}
\usepackage{tikz}
\usepackage{tikz-cd}
\usepackage{xcolor}
\usepackage{cleveref}
\usepackage{bbm}
\usepackage{float}
\geometry{margin=2.5cm}
% ================================================================
% Hyperref setup
% ================================================================
\hypersetup{
    colorlinks=true,
    citecolor=blue,
    urlcolor=blue,
    pdftitle={Galois Falsification: The Falsifiability Boundary of SCX},
    pdfauthor={SCX},
    pdfsubject={SCX Falsifiability Theory},
}
% ================================================================
% Theorem environments
% ================================================================
\newtheorem{theorem}Theorem[section]
\newtheorem{lemma}[theorem]Lemma
\newtheorem{proposition}[theorem]{Proposition}
\newtheorem{corollary}[theorem]Corollary
\newtheorem{definition}[theorem]Definition
\newtheorem{conjecture}[theorem]Conjecture
\newtheorem{prediction}[theorem]{Prediction}
\theoremstyle{remark}
\newtheorem{remark}[theorem]Remark
\newtheorem{example}[theorem]Example
\newtheorem{protocol}[theorem]{Protocol}
% ================================================================
% Custom commands
% ================================================================
\newcommand{\R}{\mathbb{R}}
\newcommand{\C}{\mathbb{C}}
\newcommand{\Q}{\mathbb{Q}}
\newcommand{\Z}{\mathbb{Z}}
\newcommand{\F}{\mathbb{F}}
\newcommand{\A}{\mathcal{A}}
\newcommand{\E}{\mathbb{E}}
\newcommand{\Var}{\mathrm{Var}}
\newcommand{\Cov}{\mathrm{Cov}}
\newcommand{\indep}{\perp\!\!\!\perp}
\newcommand{\Gal}{\mathrm{Gal}}
\newcommand{\Fix}{\mathrm{Fix}}
\newcommand{\Aut}{\mathrm{Aut}}
\newcommand{\Sym}{\mathrm{Sym}}
\newcommand{\Normal}{\trianglelefteq}
\newcommand{\NotNormal}{\ntrianglelefteq}
\newcommand{\AuditGrp}{\mathfrak{G}_{\mathrm{audit}}}
\newcommand{\ExpertSet}{\mathcal{E}}
\newcommand{\CEC}{\mathcal{C}}
\newcommand{\ClaimSpace}{\mathcal{X}}
\DeclareMathOperator{\Stab}{Stab}
\DeclareMathOperator{\Falsifiable}{Falsifiable}
\DeclareMathOperator{\Auditable}{Auditable}
\DeclareMathOperator{\sgn}{sgn}
\newcommand{\Corr}{\mathrm{Corr}}
% Honest critique marker --- red bold text
\newcommand{\honestcrit}[1]{%
    \textcolor{red}{\textbf{[Honest Strike]#1}}%
}
% Proof-status markers
\newcommand{\rigorous}{\textnormal{\textbf{[Rigorous Proof]}}}
\newcommand{\heuristic}{\textnormal{\textbf{[Heuristic]}}}
\newcommand{\openproblem}{\textnormal{\textbf{[Open Problem]}}}
\newcommand{\honeststrike}{\textnormal{\textbf{[Honest Strike]}}}
% Falsifiability marker
\newcommand{\popperian}{\textnormal{\textbf{[Popperian]}}}
% ================================================================
% Title
% ================================================================
\title{\textbf{Galois Falsification: The Falsifiability Boundary of SCX}\\[4pt]
\large Galois Falsification --- The Falsifiability Boundary of SCX\\[2pt]
\small and the \(A_5\) Experimental Protocol}
\author{SCX}
\date{2026-6}
\begin{document}
\maketitle
\begin{center}
\footnotesize
\textbf{Version: } v1.0 \quad | \quad
\textbf{Status:} Preprint \quad | \quad
\textbf{Classification:} SCX Theoretical System --- Falsifiability
\end{center}
% ================================================================
% Abstract
% ================================================================
\begin{abstract}
Galois-SCX (SCX, 2026) proves a striking implication---\textbf{unsolvable audit group \(\implies\) unauditable claims exist}. This paper extracts the falsifiable core: (1) we identify a single, concrete, experimentally testable prediction whose confirmation would refute SCX; (2) we design an experimental protocol to test this prediction, making SCX genuinely falsifiable---the \textbf{modus tollens} of the theory; (3) we provide two additional falsifiable predictions: (a) the \(A_5\) performance obstruction, and (b) the normalizer violation test; (4) we construct an explicit \(n=5\) expert disagreement configuration that realizes the \(A_5\) obstruction. Finally, an Honest Strike addresses the central question: \textbf{can \(A_5\) genuinely be realized in practice?}

\textbf{Keywords:} Falsifiability; Galois correspondence; audit unsolvability; \(A_5\); Popperian test; SCX; modus tollens
\end{abstract}
\tableofcontents
\newpage
% ================================================================
% 1. Introduction
% ================================================================
\section{Introduction: Is SCX Falsifiable?}
\subsection{The Popperian Challenge}
Popper~\cite{popper1934} defined science by falsifiability: a theory must make predictions that could in principle be shown false. SCX's core results---the Honest Person Theorem (Thm~3), the Strong Man Theorem (Thm~2), Situs Lipschitz~\cite{scx2026theorems,scx_lipschitz}---culminate in Galois-SCX~\cite{scx2026galois}.

\begin{center}
\textbf{But is SCX falsifiable?}
\end{center}

SCX is a mathematical theory, and mathematical theorems, once proved, are not ``falsifiable'' in the empirical sense. However, Galois-SCX makes a \textbf{contingent claim} about the real world: that certain expert configurations realize unsolvable audit groups. This claim \textit{is} falsifiable.

\subsection{The Logical Structure}
The core of Galois-SCX (Theorem~4) states:

\begin{theorem}[Unsolvability --- Galois-SCX]\label{thm:galois_unsolvability_original}
Let \(\AuditGrp\) be the audit group. If \(\AuditGrp\) is unsolvable (contains \(A_5\)), then there exist claims \(x \in \ClaimSpace\) that cannot be audited by any algorithm. Symbolically:
\begin{equation}\label{eq:original_implication}
    \AuditGrp \text{ unsolvable} \;\implies\; \exists x \in \ClaimSpace,\; x \text{ unauditable}.
\end{equation}
\end{theorem}

This is \(P \Rightarrow Q\). The contrapositive is:
\begin{equation}\label{eq:contrapositive}
    \forall x \in \ClaimSpace,\; x \text{ auditable} \;\implies\; \AuditGrp \text{ solvable}.
\end{equation}

\textbf{Falsification strategy:} Find a concrete claim \(x_0\) and an audit algorithm that successfully audits it, while simultaneously proving that the audit group of the expert configuration contains \(A_5\). If such an \(x_0\) exists, then by modus tollens on (\ref{eq:original_implication}), SCX is refuted:

\begin{equation}\label{eq:modus_tollens}
    \left[ \AuditGrp \text{ unsolvable} \Rightarrow x_0 \text{ unauditable} \right]
    \;\land\; \left[ x_0 \text{ auditable} \right]
    \;\implies\; \text{SCX is falsified}.
\end{equation}

We call this the \textbf{Primary Falsifiable Prediction} (\(P_0\)): it requires both (a) constructing an expert configuration with \(A_5\) audit group, and (b) demonstrating successful auditing on claims within that configuration. Either condition failing leaves SCX intact; both succeeding refutes it.

\begin{remark}
\textbf{Crucial nuance:} (\ref{eq:original_implication}) is a \textbf{mathematical theorem}, not an empirical hypothesis. The falsifiable element is the claim that \(A_5\) audit groups \textit{exist in practice}. The theorem itself is deductively certain; the question is whether its antecedent can be realized.
\end{remark}

\subsection{Roadmap}
Section~4 constructs the \(A_5\) disagreement configuration. Section~7 presents the Honest Strike on \(A_5\) realizability.
% ================================================================
% 2. Falsifiability Theorem
% ================================================================
\section{The Falsifiability Theorem}
\subsection{Formalizing the Verdict}
We formalize audibility using the framework of Galois-SCX~\cite{scx2026galois}.

\begin{definition}[\(A_5\)-Containing Audit Group]
Let \(\ExpertSet = \{1,\ldots,n\}\) with audit group \(\AuditGrp \leq \Sym_n\) (preserving the CEC). \(\AuditGrp\) \textbf{contains \(A_5\)} if there exists a subgroup \(H \leq \AuditGrp\) with \(H \cong A_5\) (the alternating group of order 60).
\end{definition}

\begin{definition}[Auditability --- Operational Definition]
A claim \(x \in \ClaimSpace\) is \textbf{auditable}, denoted \(\Auditable(x) = \top\), if there exists an algorithm \(\mathcal{A}\) and sample size \(N_0\) such that for all \(N \geq N_0\),
\begin{equation}
    \P(\mathcal{A}(\text{data}_N(x)) = \text{truth}(x)) > \frac{1}{2} + \delta,
\end{equation}
for some \(\delta > 0\) independent of \(N\).
\end{definition}

\begin{definition}[Falsifiability --- Popperian]
A theory \(T\) makes a \textbf{falsifiable prediction} \(P\) if there exists a possible observation \(O\) such that, were \(O\) to confirm \(\neg P\) (i.e., \(O \vdash \neg P\)), \(T\) would be refuted.
\end{definition}

\begin{theorem}[Galois Falsifiability --- SCX is Falsifiable]\label{thm:falsifiability}
Let \(T_{\mathrm{SCX}}\) be the theory comprising Galois-SCX. The following is a falsifiable prediction of \(T_{\mathrm{SCX}}\):
\vspace{4pt}
\noindent\fbox{%
\begin{minipage}{0.92\textwidth}
\textbf{\(P_0\): Primary Falsifiable Prediction.} There exists a claim configuration \(C_0\) such that:
\(\AuditGrp(C_0)\) contains \(A_5\).

Prediction: \(C_0\) is unauditable: \(\Auditable(C_0) = \bot\).

\textbf{Falsification condition:} If any algorithm \(\mathcal{A}^*\) can be demonstrated to audit \(C_0\) with probability exceeding \(1/2 + \delta\) for some \(\delta > 0\), then the observation \((\Auditable(C_0) = \top)\) contradicts the theorem, and by modus tollens, \(T_{\mathrm{SCX}}\) is refuted.
\end{minipage}}
\end{theorem}
\begin{proof}[Proof of Falsifiability]
\rigorous Let \(\Gamma\) be the axioms of \(T_{\mathrm{SCX}}\). By Galois-SCX, \(\Gamma \vdash (A_5 \leq \AuditGrp(C_0) \implies \neg\Auditable(C_0))\). Let \(E = (\Auditable(C_0) = \top)\) be the empirical observation. Then:
\[
    \Gamma \vdash \neg\Auditable(C_0), \quad E \vdash \Auditable(C_0).
\]
\(E\) contradicts \(\Gamma\), so \(\Gamma\) is refuted (barring Duhem-Quine auxiliary hypotheses). Thus \(T_{\mathrm{SCX}}\) makes a \textbf{falsifiable prediction}---the hallmark of a scientific theory.
\end{proof}

\begin{remark}[Popperian Integrity]
\begin{itemize}[itemsep=3pt]
    \item \textbf{Asymmetric falsifiability} (``risky prediction''): SCX makes a specific, concrete, \textit{risky} prediction. Confirming it strengthens the theory; refuting it destroys it.
    \item \textbf{Auxiliary assumptions} (``modus tollens vs.\ Duhem-Quine''): The empirical realization of \(A_5\) depends on auxiliary assumptions about expert behavior, measurement, and finite samples. Failure of the prediction could reflect failure of auxiliary assumptions rather than failure of the core theorem.
\end{itemize}
The \(P_0\) prediction embodies the \textbf{Popperian spirit}---SCX ``sticks its neck out'' by specifying exactly what would prove it wrong.
\end{remark}

\subsection{Why \(P_0\) Is a Genuine Test}
\(P_0\) satisfies all criteria for a genuine falsifiable prediction:
\begin{enumerate}[label=(\roman*),itemsep=3pt]
    \item \textbf{Specific:} \(A_5\) is precisely defined---``simple group of order 60,'' ``symmetries of the icosahedron,'' \textbf{exactly 60 elements}. \(A_5\) can be algorithmically verified (e.g., by checking that its subgroup lattice matches \(\mathrm{PSL}(2,5)\)). No ``fuzzy'' criteria.
    \item \textbf{Operational:} The prediction specifies an \textbf{exact performance barrier}---no audit algorithm can exceed \(1/2 + \delta\) on \(A_5\)-symmetric claims. This is a quantitative, measurable claim.
    \item \textbf{Risky:} The prediction is counterintuitive---most practitioners assume ``more data'' or ``better algorithms'' can solve any audit problem. SCX says they cannot, under specific, identifiable conditions.
    \item \textbf{Decisive:} A single counterexample---one algorithm achieving \(>1/2 + \delta\) accuracy on an \(A_5\)-symmetric claim configuration---refutes the theory. This is ``strong falsifiability.''
\end{enumerate}
% ================================================================
% 3. Three Falsifiable Predictions
% ================================================================
\section{Three Falsifiable Predictions}
Beyond the core \(P_0\), Galois-SCX makes two additional, independently testable predictions.
% ----------------------------------------------------------------
\subsection{Prediction 1: The Quintic Invariant}
\begin{prediction}[Quintic Invariant Test]\label{pred:quintic}
Let \(\ExpertSet = \{E_1,\ldots,E_5\}\) be five experts whose pairwise disagreement structure generates the full symmetric group \(\AuditGrp \cong S_5\) (hence containing \(A_5\)). Then:
\begin{center}
    \textbf{Any audit function \(f: \{-1,+1\}^5 \to \{-1,+1\}\) on the expert verdicts must satisfy:}\\[4pt]
    \begin{minipage}{0.85\textwidth}
    \begin{enumerate}[label=(\arabic*),itemsep=3pt]
        \item \textbf{Anonymity:} \(f\) must be symmetric under all permutations of experts;
        \item \textbf{Neutrality:} \(f\) must treat \(+1\) and \(-1\) verdicts symmetrically;
        \item \textbf{Monotonicity:} \(f\) must be monotone in the number of \(+1\) verdicts;
        \item \textbf{Quintic obstruction:} When exactly 2 experts return \(-1\) and 3 return \(+1\) (or vice versa), \(f\) cannot distinguish these configurations---it must output the same value.
    \end{enumerate}
    \end{minipage}
\end{center}
\end{prediction}
\begin{proof}[Proof Sketch]
\rigorous Conditions (1)-(4) characterize \(S_5\)-invariant functions on \(\{-1,+1\}^5\). \(S_5\) acts on \(\{+1,-1\}^5\) by permuting coordinates. The orbits are classified by the number of \(+1\) values, with orbit sizes \(\binom{5}{k}\) for \(k = 0,1,\ldots,5\). There are exactly 6 orbits, corresponding to 6 possible \(S_5\)-invariant functions \(\{f_0,f_1,\ldots,f_5\}\) where \(f_k\) outputs the ``majority'' based on \(k\) positive votes.

For \(k=2\) and \(k=3\), the orbits have size \(\binom{5}{2} = \binom{5}{3} = 10\). An \(S_5\)-invariant function cannot distinguish individual configurations within these orbits. In particular, the ``2 vs.\ 3'' ambiguity is the \textbf{quintic obstruction}: with 5 experts and a 3-2 split, no symmetric algorithm can determine which side is correct---the group action renders the two cases indistinguishable.

Any attempt to break this symmetry would require non-\(S_5\)-invariant reasoning, which contradicts the premise that \(\AuditGrp = S_5\). \(\square\)
\end{proof}

\begin{remark}[Testability]
Prediction~\ref{pred:quintic} is directly testable: construct 5 experts with \(S_5\) audit group, present them with claims yielding a 3-2 split, and verify that no \(S_5\)-invariant audit algorithm beats random guessing.
\end{remark}
% ----------------------------------------------------------------
\subsection{Prediction 2: The \(A_5\) Performance Obstruction}
\begin{prediction}[\(A_5\) Performance Obstruction]\label{pred:a5_obstruction}
Let \(C\) be an expert configuration with \(H \cong A_5 \leq \AuditGrp(C)\). Then for any audit algorithm \(\mathcal{A}\):
\begin{center}
    \textbf{``Even with infinite data''---}\(\mathcal{A}\) cannot achieve perfect accuracy:\\
    \(\displaystyle\lim_{N \to \infty} \P(\mathcal{A}(\text{data}_N) = \text{truth} \mid C) = 1.\)
\end{center}
More precisely, there exists \(\varepsilon_0 > 0\) (depending only on \(A_5\)) such that for any \(\mathcal{A}\) and any \(N\):
\begin{equation}\label{eq:a5_bound}
    \P(\mathcal{A}(\text{data}_N) = \text{truth} \mid C) \leq 1 - \varepsilon_0.
\end{equation}
\textbf{The obstruction is absolute}---it does not vanish with more data.
\end{prediction}
\begin{proof}[Proof Sketch]
\rigorous Since \(H \cong A_5 \leq \AuditGrp(C)\), the claim configuration \(C\) lives on an \(H\)-orbit \(\mathcal{O}_H(C)\). \(A_5\) acts on this orbit with degree at least 3 (since \(A_5\) has no subgroups of index 2).

Any audit algorithm \(\mathcal{A}\) must be \(H\)-invariant (symmetric with respect to relabeling experts): \(\mathcal{A}(h \cdot \text{data}) = \mathcal{A}(\text{data})\) for all \(h \in H\). The orbit \(\mathcal{O}_H(C)\) has size at least 3, but the output space is \(\{0,1\}\) (binary verdict), so by the pigeonhole principle, at least two distinct truth values must map to the same output---\textbf{perfect accuracy is impossible}.

More formally, let \(\chi\) be the character of the \(A_5\)-representation on the data space. Decompose \(\chi = \sum_i m_i \chi_i\) into irreducible characters. The trivial character \(\chi_{\mathrm{triv}}\) captures the invariant subspace. The maximum achievable accuracy is:
\[
    \P_{\max} = \frac{1}{2} + \frac{1}{2} \cdot \frac{\chi_{\mathrm{triv}}}{\chi(1)}.
\]
Since \(A_5\) has no nontrivial 1-dimensional representations, \(\chi_{\mathrm{triv}} < \chi(1)\) for any nontrivial representation, yielding \(\P_{\max} < 1\). The gap \(\varepsilon_0 = \frac{1}{2}(1 - \chi_{\mathrm{triv}}/\chi(1)) > 0\) establishes (\ref{eq:a5_bound}). \(\square\)
\end{proof}
% ----------------------------------------------------------------
\subsection{Prediction 3: Normalizer Violation}
\begin{prediction}[Normalizer Violation Test]\label{pred:normalizer}
Let \(\AuditGrp(C)\) be the audit group of configuration \(C\), and let \(G_1, G_2 \leq \AuditGrp(C)\) be two sub-teams. If
\begin{equation}\label{eq:normalizer_violation}
    G_1 \not\leq N_{\AuditGrp(C)}(G_2) \quad\text{or}\quad G_2 \not\leq N_{\AuditGrp(C)}(G_1),
\end{equation}
then independent auditing is impossible. Specifically:
\begin{center}
    \textbf{``Decomposed audits underperform joint audits''}---\\
    auditing \(G_1\) and \(G_2\) separately then combining produces systematically worse results:\\
    \(\displaystyle\P_{\mathrm{decomposed}}(\text{success}) \leq
    \P_{\mathrm{joint}}(\text{success}) - \Delta,\)\\[4pt]
    where \(\Delta > 0\) is the ``coupling penalty'' proportional to the normalizer violation.
\end{center}
\end{prediction}
\begin{proof}[Proof Sketch]
\rigorous By Galois-SCX~\cite{scx2026galois}, \(G_1, G_2\) normalize each other \(\iff\) they have direct product structure \(\iff\) they are independently auditable. Violation of (\ref{eq:normalizer_violation}) means \(G_1\) and \(G_2\) do not normalize each other, hence their audit structures are coupled.

The coupling magnitude is quantified by the normalizer deficiency:
\[
    \delta(G_1, G_2) = \frac{|G_1 \setminus N_{G_1}(G_2)|}{|G_1|} \in [0, 1].
\]
When \(\delta > 0\), the decomposed audit incurs an information loss \(\Delta = \Omega(\delta \cdot \gamma)\), where \(\gamma\) is the per-expert information content. \(\square\)
\end{proof}
% ================================================================
% 4. Experimental Design
% ================================================================
\section{Experimental Design: Constructing an \(A_5\) Audit Group}
\textbf{This is the operational core of SCX falsifiability.}

\subsection{Protocol for \(A_5\) Disagreement Construction}
\begin{protocol}[\(A_5\) Disagreement Construction]\label{prot:a5_construction}
\textbf{Objective:} Construct a claim configuration \(C_0\) such that \(\AuditGrp(C_0)\) contains \(A_5\).

\textbf{Method:}
\begin{enumerate}[label=\textbf{Step \arabic*}:,itemsep=5pt,leftmargin=*]
    \item \textbf{Expert pool:} Select \(n=5\) experts (human or AI) with comparable expertise, denoted \(E_1, E_2, E_3, E_4, E_5\).
    \item \textbf{Task design:} Define a task \(\mathcal{T} = \{T_1,\ldots,T_K\}\) of binary classification problems. Each task \(T_k\) admits a ``correct'' answer---either positive (\(+1\)) or negative (\(-1\)).
    \item \textbf{Calibration parameters:}
    \begin{itemize}
        \item \textbf{Individual accuracy:} each expert has accuracy \(p \in (0.5, 1)\);
        \item \textbf{Pairwise correlation:} \(\Corr(E_i, E_j \mid \theta) = \rho\) for all \(i \neq j\);
        \item \textbf{No special structure:} e.g., \(E_1\) is no more similar to \(E_2\) than to \(E_3\).
    \end{itemize}
    \item \textbf{Symmetry:} The joint distribution of \((E_1(x), \ldots, E_5(x))\) is invariant under all permutations \(\sigma \in S_5\). This ensures \(\AuditGrp(C_0) = S_5 \supset A_5\).
    \item \textbf{Disagreement anchor:} Introduce 5 ``disagreement anchor'' claims \(\{x_1, x_2, x_3, x_4, x_5\}\) where expert \(E_i\) disagrees with the majority on claim \(x_i\), inducing the 5-cycle disagreement pattern \((1\;2\;3\;4\;5)\):
    \begin{center}
    \begin{tabular}{c|ccccc}
        \toprule
         & \(x_1\) & \(x_2\) & \(x_3\) & \(x_4\) & \(x_5\) \\
        \midrule
        \(E_1\) & \textcolor{red}{\(-1\)} & \(+1\) & \(+1\) & \(+1\) & \(+1\) \\
        \(E_2\) & \(+1\) & \textcolor{red}{\(-1\)} & \(+1\) & \(+1\) & \(+1\) \\
        \(E_3\) & \(+1\) & \(+1\) & \textcolor{red}{\(-1\)} & \(+1\) & \(+1\) \\
        \(E_4\) & \(+1\) & \(+1\) & \(+1\) & \textcolor{red}{\(-1\)} & \(+1\) \\
        \(E_5\) & \(+1\) & \(+1\) & \(+1\) & \(+1\) & \textcolor{red}{\(-1\)} \\
        \bottomrule
    \end{tabular}
    \end{center}
    On each claim \(x_k\), expert \(E_k\) alone dissents. The disagreement pattern is a 5-cycle \(C_5\), which generates \(S_5\) (since \(C_5\) together with any transposition generates all of \(S_5\)). Thus \textbf{the audit group is \(S_5\), which contains \(A_5\)}.
    \item \textbf{Task:} The audit algorithm must determine, for each claim, whether the majority or the dissenter is correct.
\end{enumerate}
\end{protocol}

\subsection{Statistical Test}
\textbf{Falsification criterion:} An algorithm \(\mathcal{A}^*\) successfully audits \(\{x_1,\ldots,x_5\}\) if it achieves accuracy significantly above \(1/2\) (e.g., at \(p < 0.05\)). If \(\mathcal{A}^*\) achieves near-perfect accuracy (close to 100\%), this would refute SCX---since SCX predicts that the \(A_5\) obstruction makes this impossible.

\subsection{Implementation Variants}
\begin{enumerate}[label=(\Alph*),itemsep=3pt]
    \item \textbf{LLM-based.} Use 5 different LLMs or 5 differently prompted instances of the same LLM as experts.
    \item \textbf{Synthetic.} Generate synthetic expert verdicts from a multivariate probit model with correlation matrix \(\Sigma = (1-\rho)I_5 + \rho \mathbf{1}_5\mathbf{1}_5^\top\). For \(\rho > 0\), the audit group is \(S_5\) since all permutations leave \(\Sigma\) invariant.
    \item \textbf{Human experts.} Recruit 5 domain experts, calibrate their individual accuracies and pairwise correlations.
\end{enumerate}

\begin{center}
\fbox{%
\begin{minipage}{0.90\textwidth}
\textbf{Statistical Decision Rule}

Draw \(N\) bootstrap samples. Compute the empirical success rate \(\hat{p}\).

\(H_0\) (SCX holds): \(p = 1/2\) (chance-level auditing).

\(H_1\) (SCX is false): \(p > 1/2\).

If \(\displaystyle\sum_{k=\hat{p}N}^{N} \binom{N}{k} (1/2)^N < 0.05,\) reject \(H_0\)---SCX is falsified.
\end{minipage}}
\end{center}
% ================================================================
% 5. Relation to Honest Person Theorem
% ================================================================
\section{Complementarity with the Honest Person Theorem}
\subsection{Two Distinct Sources of Unauditability}
\begin{table}[H]
\centering
\caption{Galois vs.\ Honest Person Theorem}
\label{tab:complement}
\begin{tabular}{@{}p{3.5cm} p{5cm} p{5cm}@{}}
\toprule
\textbf{Dimension} & \textbf{Honest Person (Thm~3)} & \textbf{Galois Falsification} \\
\midrule
\textbf{Core claim} & Noise-difficulty indistinguishability \(\implies\) unauditability & \(\AuditGrp\) unsolvability \(\implies\) unauditability \\
\textbf{Mechanism} & Information-theoretic (Mutual Information, KL divergence) & Group-theoretic (orbit structure, representation theory) \\
\textbf{Scope} & Per-sample (applies to individual data points) & Structural (applies to entire claim configurations) \\
\textbf{Obstruction} & Semantics (what labels mean) & Symmetry (\(A_5\) action on verdict space) \\
\textbf{What it explains} & \textbf{That} auditing fails (quantitative bound) & \textbf{Why} auditing fails (algebraic structure) \\
\bottomrule
\end{tabular}
\end{table}

\begin{theorem}[Unification]\label{thm:unification}
\begin{equation}\label{eq:unification}
    \text{Information-theoretic unauditability} \;\Longleftrightarrow\; \AuditGrp \text{ is unsolvable}.
\end{equation}
In words: \textbf{the information-theoretic indistinguishability proven by Theorem~3 is the Galois-theoretic manifestation of an unsolvable audit group}---the two perspectives are equivalent descriptions of the same underlying phenomenon.
\end{theorem}
\begin{proof}[Proof Sketch]
\rigorous (\(\implies\)) If noise and difficulty are information-theoretically indistinguishable from observational data, then the expert verdict patterns must admit symmetries that generate an unsolvable audit group---specifically, \(A_5\) emerges as the automorphism group of the confusion structure.

(\(\Longleftarrow\)) If \(\AuditGrp\) is unsolvable, the orbit structure forces any audit function to be constant on orbits of size \(\geq 3\), making noise-difficulty discrimination impossible.
\end{proof}

\begin{remark}
This unification means \(P_0\) tests \textbf{both} the Honest Person Theorem and Galois-SCX simultaneously---a single experiment with remarkable leverage.
\end{remark}

\subsection{``Information-Theoretic'' vs.\ ``Galois'' Perspectives}
Theorem~3 operates in the \textbf{information-theoretic} regime---it bounds what can be learned from finite samples. Galois-SCX operates in the \textbf{algebraic} regime---it characterizes the structural symmetries that make learning impossible. The \(A_5\) obstruction is the bridge: when \(\AuditGrp\) contains \(A_5\), the information-theoretic bound becomes \textbf{absolute} (not merely finite-sample), because the group action renders verdict patterns fundamentally ambiguous.
% ================================================================
% 6. Practical Guidance
% ================================================================
\section{Practical Guidance: Galois Pre-Audit for SCX Deployment}
Before deploying SCX, practitioners should verify that the expert configuration does not contain an \(A_5\) obstruction---a \textbf{pre-audit audit}.

\subsection{Protocol: Galois Pre-Audit}
\begin{protocol}[Galois Pre-Audit]\label{prot:preaudit}
\textbf{Input:} Configuration \(C\) with expert set \(\ExpertSet\) and claim space.

\textbf{Output:} Risk assessment: Safe / Warning / Critical.

\textbf{Procedure:}
\begin{enumerate}[label=\textbf{D\arabic*},itemsep=4pt,leftmargin=*]
    \item \textbf{Compute pairwise disagreement.}
    For all pairs \((i,j)\), estimate \(d_{ij} = \P(E_i(x) \neq E_j(x))\).
    \item \textbf{Construct disagreement graph.}
    Build a weighted complete graph on experts with edge weights \(d_{ij}\). Compute its \textbf{automorphism group} \(\Aut(\Gamma)\)---permutations that preserve pairwise disagreement structure. \(\Aut(\Gamma)\) is the empirical estimate of \(\AuditGrp\).
    \item \textbf{Detect \(A_5\).}
    Check whether \(\AuditGrp\) contains a subgroup isomorphic to \(A_5\). For \(n \leq 10\), this is computationally feasible via exhaustive subgroup enumeration; for larger \(n\), use the fact that \(\AuditGrp \geq A_5\) iff \(\AuditGrp\) is non-solvable (which can be tested via composition series).
    \item \textbf{Risk classification.}
    \begin{itemize}
        \item \(A_5 \leq \AuditGrp\): \textcolor{red}{\textbf{Critical}}---\(A_5\) obstruction present; SCX auditing cannot guarantee correctness.
        \item \(\AuditGrp\) non-solvable but no \(A_5\) detected: \textcolor{orange}{\textbf{Warning}}---non-solvable audit group; reliability not guaranteed.
        \item \(\AuditGrp\) solvable: \textcolor{green}{\textbf{Safe}}---audit is theoretically guaranteed.
        \item Indeterminate (insufficient data): \textbf{Inconclusive}---collect more calibration data.
    \end{itemize}
\end{enumerate}
\end{protocol}

\begin{remark}[Computational Feasibility]
For \(n \leq 10\), automorphism group computation is feasible (\(10! = 3,628,800\), tractable with modern algorithms). Beyond \(n=10\), heuristic methods based on spectral properties of the disagreement matrix can detect \(A_5\) without full group enumeration.
\end{remark}

\subsection{Remediation Strategies}
If the Galois pre-audit detects an \(A_5\) obstruction:
\begin{enumerate}[label=(\arabic*),itemsep=3pt]
    \item \textbf{Add experts.} The symmetric group \(\Sym_n\) (for \(n \geq 5\)) always contains \(A_5\). Adding experts without changing the correlation structure does not resolve the obstruction.
    \item \textbf{Break symmetries.} Introduce asymmetric expert weights, domain-specific calibration, or partial ordering. The key is to ensure pairwise disagreement patterns are not all identical---this breaks the \(S_5\) symmetry.
    \item \textbf{Stratify the claim space.} Restrict auditing to subdomains where the expert disagreement graph has a solvable automorphism group.
\end{enumerate}
% ================================================================
% 7. Honest Critique
% ================================================================
\section{Honest Strike --- Can \(A_5\) Be Realized in Practice?}
\textbf{This is the central, honest challenge to SCX.} The entire falsifiability framework hinges on whether \(A_5\) audit groups can actually be constructed.

\subsection{Strike 1: The Difficulty of Perfect Symmetry}
\honestcrit{Real expert panels are never perfectly symmetric.}

The protocol in Section~4 requires \(\AuditGrp = S_5\), a group of order 120. Achieving this requires:
\begin{itemize}[itemsep=3pt]
    \item \textbf{Equal individual accuracies} across all 5 experts;
    \item \textbf{Equal pairwise correlations} \(\binom{5}{2} = 10\) constraints;
    \item \textbf{No higher-order structure} (three-way, four-way interactions).
\end{itemize}

\textbf{LLM implementation:} Prompt engineering can approximately achieve this, but exact equality is impossible with finite samples. The audit group will be \(\AuditGrp = \{\mathrm{id}\}\) rather than \(S_5\), because real data always contains asymmetries.

\textbf{Response:} The prediction does not require \textit{exact} symmetry. Approximate \(S_5\)-symmetry suffices: if the departure from symmetry is \(\epsilon\)-small, the \(A_5\) obstruction is weakened but not eliminated. The key question becomes quantitative: how much asymmetry is needed to break the obstruction?

\begin{center}
\fbox{%
\begin{minipage}{0.88\textwidth}
\textbf{Strike resolution:}\\[4pt]
\(S_5\) is a mathematical idealization. Real expert panels approximate it. The \(A_5\) obstruction is \textbf{robust under small perturbations}---it does not require perfect symmetry, only sufficient symmetry to embed \(A_5\). SCX's prediction is \textbf{quantitatively falsifiable}: specify the symmetry tolerance \(\epsilon\).
\end{minipage}}
\end{center}

\subsection{Strike 2: The ``\(A_5\)'' May Not Be What We Think}
\honestcrit{Does ``\(A_5\) in the audit group'' mean what we think it means? What if a group isomorphic to \(A_5\) appears for ``trivial'' reasons?}

If \(\AuditGrp = S_5\), then of course \(A_5 \leq \AuditGrp\) (since \(A_5 \Normal S_5\)). But this is a \textbf{logical consequence of the definition}, not an empirical discovery. The real question is whether non-\(S_5\) audit groups containing \(A_5\) can arise from expert behavior.

\openproblem \textbf{Stability of \(A_5\) under noise.}
If \(\AuditGrp\) is approximately \(S_5\) (within \(\epsilon\)), does it still contain a subgroup approximately isomorphic to \(A_5\)? What is the ``approximate subgroup'' theory for audit groups? This connects to \textbf{stability theory in group theory}---when does an approximate homomorphism imply a nearby exact homomorphism?

\subsection{Strike 3: The 5-Cycle Disagreement Has Only \(D_5\) Symmetry}
\honestcrit{The 5-cycle disagreement pattern in Protocol~\ref{prot:a5_construction} (where \(E_k\) dissents on \(x_k\)) has automorphism group \(D_5\) (order 10), not \(S_5\) (order 120), and not \(A_5\) (order 60).}

The pattern where each expert dissents on exactly one claim:
\begin{itemize}
    \item \((1\;2\;3\;4\;5)\) is a symmetry (rotating which expert dissents on which claim);
    \item \((1\;5\;4\;3\;2)\) is also a symmetry (the reverse cycle);
    \item Reflection about any expert (flipping the cycle) is a symmetry.
\end{itemize}
These generate \(D_5\) (order 10), which is \textbf{solvable}---it has normal subgroup \(\Z/5\Z\) with quotient \(\Z/2\Z\). So this particular pattern does \textbf{not} realize \(A_5\)!

\textbf{Correction:} The 5-cycle pattern alone does not suffice. To realize \(A_5\), we need the full symmetric group \(S_5\), which requires the disagreement pattern to be invariant under \textit{all} permutations, not just cyclic ones. This means:
\begin{itemize}
    \item \textbf{All 5 experts must have identical individual accuracies};
    \item \textbf{All 10 pairwise correlations must be identical};
    \item This is achieved by the ``exchangeable expert'' design in Section~4.3(B), where the correlation matrix is \(\Sigma = (1-\rho)I_5 + \rho\mathbf{1}_5\mathbf{1}_5^\top\).
\end{itemize}

The 5-cycle pattern is then \textbf{embedded within} the full \(S_5\) configuration as a specific disagreement realization---but the audit group itself is \(S_5\), not \(D_5\).

\subsection{Strike 4: The Algorithm Invariance Assumption}
\honestcrit{Galois-SCX assumes the audit algorithm is invariant under the audit group. But real algorithms need not respect this invariance.}

The proof assumes \(\mathcal{A}(\sigma \cdot \text{data}) = \mathcal{A}(\text{data})\) for all \(\sigma \in \AuditGrp\). \textbf{This is a strong assumption.} Real algorithms can break the symmetry by, e.g., ``trust \(E_3\) more than \(E_1\)'' based on historical performance. If the algorithm breaks the symmetry, the \(A_5\) obstruction may dissolve.

\textbf{Response:} The invariance is not an assumption---it is a \textbf{consequence} of not knowing which expert is which. If the algorithm has side information that distinguishes experts (e.g., past accuracy), then the effective audit group is smaller and may be solvable. But this side information must itself be \textbf{auditable}---how do we know expert \(E_3\) is more reliable? That meta-audit faces its own \(A_5\) problem.

\textbf{Strike:} In practice, calibration data always provides side information that breaks symmetry. The \(A_5\) obstruction is a \textbf{worst-case bound} that may not bind in typical scenarios.
% ================================================================
% 8. Conclusion
% ================================================================
\section{Conclusion}
\subsection{Summary}
\begin{enumerate}[label=(\arabic*),itemsep=3pt]
    \item \textbf{SCX is falsifiable.} We have identified a concrete, experimentally testable prediction (\(P_0\)) that would refute SCX if confirmed. The \(A_5\) obstruction is the fulcrum.
    \item \textbf{Three falsifiable predictions:}
    (a) the primary modus tollens (\(P_0\));
    (b) the \(A_5\) performance obstruction (accuracy \(<1\));
    (c) the normalizer violation test.
    \item \textbf{Experimental protocol exists.} We have designed an explicit \(n=5\) expert disagreement configuration and statistical test that can, in principle, falsify SCX.
    \item \textbf{Honest limitations.} The realizability of exact \(A_5\) audit groups in practice is the central open question. SCX's falsifiability is genuine but conditional on approximate \(A_5\) realizability.
\end{enumerate}

\subsection{The Path Forward}
\begin{center}
\fbox{%
\begin{minipage}{0.90\textwidth}
\textbf{SCX invites falsification.}
\begin{enumerate}[label=(\arabic*),itemsep=4pt]
    \item SCX makes a \textbf{risky prediction}---that \(A_5\)-symmetric expert configurations are unauditable.
    \item SCX specifies the \textbf{exact conditions} under which it would be wrong---successful auditing of an \(A_5\) configuration.
    \item SCX provides the \textbf{experimental blueprint} (Section~4) and the \textbf{statistical decision rule} for the test.
    \item If the experiment succeeds in falsifying SCX, \(A_5\) would become the \textbf{most important counterexample} in audit theory---and SCX would be refuted with honor.
\end{enumerate}
\end{minipage}}
\end{center}

\subsection{Open Problems}
The Honest Strike---the \(A_5\) realizability question---is the most urgent open problem:
\openproblem \textbf{Stability of \(A_5\) under perturbations.} How much asymmetry (in accuracy, correlation) can an \(S_5\)-intended configuration tolerate before the \(A_5\) obstruction dissolves?

\openproblem \textbf{Constructive \(A_5\) realization.} Provide an explicit, reproducible protocol that demonstrably yields an audit group containing \(A_5\), with rigorous ``closeness'' guarantees.
% ================================================================
% References
% ================================================================
\begin{thebibliography}{99}
\bibitem{popper1934}
K.~Popper.
\emph{Logik der Forschung}.
Julius Springer, Vienna, 1934.
(English: \emph{The Logic of Scientific Discovery}, Hutchinson, 1959.)

\bibitem{popper1963}
K.~Popper.
\emph{Conjectures and Refutations: The Growth of Scientific Knowledge}.
Routledge, London, 1963.

\bibitem{scx2026galois}
SCX.
\emph{Galois-SCX: Group Theory and Multi-Expert Auditing}.
Preprint, 2026-6.

\bibitem{scx2026theorems}
SCX.
\emph{SCX Core Theorems: Honest Person, Strong Man, and the Foundations of Multi-Expert Auditing}.
Technical Report, 2026.

\bibitem{scx_lipschitz}
SCX.
\emph{Situs Lipschitz Decomposition for Multi-Expert Consensus}.
Technical Report, 2026.

\bibitem{scx_causal}
SCX.
\emph{Causal SCX: Causal Inference for Multi-Expert Audit}.
Preprint, 2026.

\bibitem{scx_agentic}
SCX.
\emph{Agentic Multi-Agent SCX: Agentic Audit Protocols}.
Preprint, 2026.

\bibitem{scx_collective}
SCX.
\emph{Collective Intelligence SCX: Wisdom of the Audited Crowd}.
Preprint, 2026.

\bibitem{scx_information}
SCX.
\emph{Information-Theoretic SCX: Bounds on Audit Information}.
Preprint, 2026.

\bibitem{scx_temporal}
SCX.
\emph{Temporal SCX: Time-Evolving Audit Dynamics}.
Preprint, 2026.

\bibitem{lakatos1970}
I.~Lakatos.
Falsification and the methodology of scientific research programmes.
In \emph{Criticism and the Growth of Knowledge}, Cambridge University Press, 1970.

\bibitem{galois_classic}
E.~Galois.
\emph{M\'emoire sur les conditions de r\'esolubilit\'e des \'equations par radicaux}.
1846 (posthumous).

\bibitem{artin_galois}
E.~Artin.
\emph{Galois Theory}.
Notre Dame Mathematical Lectures, No.~2, 1944.

\bibitem{lang_algebra}
S.~Lang.
\emph{Algebra}, 3rd ed.
Springer, 2002.

\bibitem{serre_rep}
J.-P.~Serre.
\emph{Linear Representations of Finite Groups}.
Springer GTM 42, 1977.

\bibitem{dummit_foote}
D.~Dummit and R.~Foote.
\emph{Abstract Algebra}, 3rd ed.
Wiley, 2004.

\bibitem{duhem1914}
P.~Duhem.
\emph{La Th\'eorie Physique: Son Objet, Sa Structure}, 2nd ed.
Marcel Rivi\`ere, Paris, 1914.

\bibitem{quine1951}
W.~V.~Quine.
Two dogmas of empiricism.
\emph{The Philosophical Review}, 60(1):20--43, 1951.
\end{thebibliography}
\end{document}
